\section{A certificate-defined series and explicit generation tails}
\label{rr:section}

The existing full-model theorems supply two complete strict searches at
a positive rational test point $q$: a finite lower-history witness for
$q<\mu$, and a finite-window upper witness for $\mu<q$. Neither search
is required to halt if $q=\mu$. The following scheduling refinement
makes the resulting rational series completely explicit while avoiding
an equality oracle. It does not supply a running-time bound for finding
the individual certificates.

\begin{proposition}[Equality-safe trisection]\label{rr:trisection}
Start from a strict rational bracket $a_0<\mu<b_0$, for example
$a_0=1/2$, $b_0=3/5$ with common denominator $Q=10$. There is a
deterministically specified sequence of strict rational brackets
$a_n<\mu<b_n$ satisfying
\[
 a_n\le a_{n+1}<b_{n+1}\le b_n,
 \qquad b_n-a_n\le(b_0-a_0)(2/3)^n.
\]
If $Q$ is a common denominator of $a_0,b_0$, then both endpoint
denominators divide $Q3^n$. In particular,
\begin{equation}\label{rr:series}
 \mu=a_0+\sum_{n\ge1}(a_n-a_{n-1}),\qquad
 0<\mu-a_n<(b_0-a_0)(2/3)^n.
\end{equation}
The summands are nonnegative rational numbers, and can vanish. The
midpoints $q_n=(a_n+b_n)/2$ have denominators dividing $2Q3^n$ and
\[
 |q_n-\mu|<\tfrac12(b_0-a_0)(2/3)^n.
\]
\end{proposition}

\begin{proof}
Fix computable enumerations of all finite lower-certificate attempts
and all finite upper-certificate attempts, including every truncation
parameter and rational enclosure precision. Each attempt is a finite
calculation. Run these enumerations as resumable algorithms, reporting
only after a strict certificate has been verified.

At stage $n$, put
\[
 x_n=(2a_n+b_n)/3,\qquad y_n=(a_n+2b_n)/3.
\]
Run the four searches in the fixed cyclic order: lower at $x_n$,
upper at $x_n$, lower at $y_n$, upper at $y_n$. At each cycle perform
one computation step of each search, and stop at the first verified
certificate in this order. This is a deterministic fair schedule.
Since $x_n\ne y_n$, at least one test point differs from $\mu$; the
corresponding strict comparison has a complete search and eventually
succeeds. Thus every stage terminates.

If the first success proves $q<\mu$, replace the lower endpoint by
$q$; if it proves $\mu<q$, replace the upper endpoint by $q$. The
new interval is strict and has width at most two thirds of the old
width. Induction gives the width estimate and shows that its endpoints
lie on the grid $(Q3^n)^{-1}\mathbb Z$. The widths tend to zero, so
both endpoints tend to $\mu$. Telescoping the lower endpoints proves
\eqref{rr:series}; the midpoint claims follow immediately.
\end{proof}

For the original exponent one may use the known strict bracket
$a_0=1/2$, $b_0=3/5$, so $Q=10$. The midpoint error is then at most
$\tfrac1{20}(2/3)^n$, with midpoint denominator dividing $20\cdot3^n$.
Using this wider initial interval avoids assuming that either endpoint
of the sharper closed enclosure differs from $\mu$.

\paragraph{Computational and arithmetic scope.}
The convergence rate counts completed certificate stages, not elementary
operations. The certificate searches can be extremely expensive. This
is a more concrete implementation of the existing computability
theorem, rather than a finite special-function expression for the
exponent or an efficient high-precision solver.

Ordering gives a genuine but weak noncancellation statement. The integer
polynomial
\[
 P_n(X)=Q3^nX-Q3^na_n
\]
satisfies $P_n(\mu)>0$, while its available error bound is only
\[
 0<P_n(\mu)<Q(b_0-a_0)2^n.
\]
It does not tend to zero after clearing denominators. The construction
also works for rational roots, including roots equal to intermediate
test points. Its nonvanishing therefore supplies no irrationality or
transcendence implication.

\begin{proposition}[Explicit geometric generation tails]
\label{rr:generation-tail}
Suppose $b\ge\mu$ is a certified positive upper bound. Let
$b<s\le S$, where $S>1$, and put
\[
 C_b=12(b+3),\qquad
 \theta=\frac{S-s}{S-b},\qquad
 \rho=S^{-(1-\theta)}<1.
\]
Then, for every $m\ge1$,
\begin{equation}\label{rr:holder-bound}
 A_m^\infty(s)\le C_b^{\theta}\rho^m.
\end{equation}
Consequently the full generation resolvent
$\mathcal Z(s)=1+\sum_{m\ge1}A_m^\infty(s)$ has the explicit tail
\begin{equation}\label{rr:generation-tail-bound}
 0\le\mathcal Z(s)-1-\sum_{m=1}^M A_m^\infty(s)
 \le\frac{C_b^{\theta}\rho^{M+1}}{1-\rho}.
\end{equation}
\end{proposition}

\begin{proof}
The existing pressure supremum formula implies
$A_m^\infty(b)\le12(b+3)r_\infty(b)^m\le C_b$, since $b\ge\mu$.
The potential genealogy gives $A_m^\infty(S)\le S^{-m}$.
Each $A_m^\infty$ is the Laplace transform of a positive measure, and
$s=\theta b+(1-\theta)S$. H\"older's inequality therefore gives
\[
 A_m^\infty(s)
 \le A_m^\infty(b)^\theta A_m^\infty(S)^{1-\theta}
 \le C_b^\theta S^{-(1-\theta)m}.
\]
At $s=S$ this is just the potential-genealogy bound. Summing the
geometric majorant proves the tail estimate.
\end{proof}

For rational $b,s,S$, no algebraic root calculation is needed to obtain
a rational geometric majorant. Put $\alpha=(s-b)/(S-b)\in(0,1]$.
Convexity of $x\mapsto S^{-x}$ on $[0,1]$ gives
\[
 \rho=S^{-\alpha}\le q:=1-\alpha(1-1/S)<1.
\]
Since $C_b>1$, a completely rational tail bound is therefore
\[
 \frac{C_bq^{M+1}}{1-q}.
\]
For $S=2$, this uses $q=1-(s-b)/[2(2-b)]$.
This gives an explicit stopping rule for the generation truncation
above any certified upper bound. At any fixed rational $s>\mu$, one
can first find a rational upper certificate $\mu<b<s$ and then use
this construction. The search for $b$ has no asserted uniform
complexity bound.

\paragraph{Why this is not a critical-pole theorem.}
The geometric ratio tends to one as $s\downarrow b$. Moreover,
$r_\infty(\mu)=1$ and bounded critical coefficients do not by
themselves imply divergence of $\sum_m A_m^\infty(\mu)$.
For example, the abstract sequence $a_m=(m+1)^{-2}$ is
supermultiplicative, is bounded, and has $a_m^{1/m}\to1$, but is
summable. Thus an assertion that the reciprocal resolvent vanishes at
the critical point still requires a separate critical nonsummability
argument.

\subsection{A sharper population tail for fixed-time calculations}

The following exact Yule identity already appears in the project's
earlier exact-characterization work. It can replace the cruder
second-moment cutoff bound in the present finite-time calculation of
$H(T)$.

For the Yule population $Y_n(T)$ from $n$ founders, put $p=e^{-T}$.
For $N\ge n$,
\begin{equation}\label{rr:yule-tail}
 \mathbb E[Y_n(T);Y_n(T)>N]
 =\frac np\sum_{j=0}^n\binom{N+1}{j}p^j(1-p)^{N+1-j}
 =\frac np I_{1-p}(N-n+1,n+1).
\end{equation}
Here $I_x$ is the regularized incomplete beta function; its finite
binomial and hypergeometric representations are standard\cite{DLMF}.
Indeed,
\[
 k\binom{k-1}{n-1}p^n(1-p)^{k-n}
 =\frac np\binom{k}{n}p^{n+1}(1-p)^{k-n}.
\]
Thus size bias replaces $Y_n$ by $Y_{n+1}-1$; its tail is the
binomial event with at most $n$ successes in $N+1$ trials. For fixed
$n,T$, the error in~\eqref{rr:yule-tail} is
$O_{n,T}(N^n(1-e^{-T})^N)$.

If $0<p_-\le e^{-T}$ is rational, the same finite expression evaluated
at $p_-$ is a rational upper bound. This follows by coupling with a
Yule process at the later time $-\log p_-$ and using the monotonicity
of $y\mathbf1_{\{y>N\}}$. The bound is uniform for $0\le t\le T$.
In a finite collection of retained initial profiles, take the maximum
over their initial populations and multiply by $T$ to bound the
population-cutoff contribution to the integral defining $H(T)$.
The spatial-profile and quadrature errors remain separate. This
geometric convergence for a fixed-time finite-state calculation does
not provide a convergence rate for $T^{-1}\log H(T)$ to $\mu$.
